\chapter{Каноническая мера и локализация лептонного разрыва}

Гипотеза H-M предлагает принять одну каноническую кинетическую меру,
сохранить нейтринную норму \(23+\pi^{-1}\), читать лептонный тангент с
нормой \(8/3\), а сырой seed \(\pi^2+2\pi+2/3\) перенести из кинетического
члена в вершину. Эта схема согласована как разложение ролей, но необходимо
различать совместимость меры, вывод меры и физический двухсекторный тест.

\section{Тест совместимости H-M}

Единичный вес \(W=I\) следует понимать относительно уже фиксированной
градуированной trace--Hodge метрики. Он не означает евклидову норму один
для каждого компонента: интегральная одна форма сохраняет геометрическую
норму \(\pi^{-1}\). При таком понимании
\begin{align}
 \|\boldsymbol\Xi_\nu\|^2
 &=23\|\widehat e_0\|^2+\|e_1\|^2
 =23+\pi^{-1},\\
 \|\boldsymbol\Xi_\tau\|^2
 &=1+1+\frac23=\frac83.
\end{align}
Следовательно, H-M проходит алгебраический тест совместимости.

Однако каноническая метрика в формулировке H-M зафиксирована как исходное
условие. Ни симметрия, ни гессиан родительского действия, ни квантовая мера,
единственным образом выбирающие её, не предъявлены. Поэтому корректный
статус равен «каноническая мера допустима», а не «каноническая мера
выведена».

\section{Физический двухсекторный гейт}

Число \(8/3\) ранее возникло как контроль канонической нормировки, которая
удаляет объёмные множители из лептонного seed. Оно не является независимой
наблюдаемой заряженного лептонного сектора. Поэтому пара
\begin{equation}
 23+\pi^{-1},
 \qquad
 \frac83
\end{equation}
задаёт два согласованных математических чтения одной метрики, но не два
закрытых физических предсказания. Число прошедших независимых
normalization-sensitive секторов остаётся равным одному.

\begin{nogocriterion}[Совместимость не является выводом меры]
Назначение \(W=I\) проходит H-M-гейт только как фиксированная кинематическая
конвенция. Для статуса выведенной меры требуется получить ту же
градуированную метрику из вариационного функционала и вычислить с её помощью
независимую наблюдаемую, не использованную при выборе нормировки.
\end{nogocriterion}

\section{Перенос сырого seed в вершину}

Исключение
\begin{equation}
 \rho_\tau^{\mathrm{raw}}
 =\pi^2+2\pi+\frac23
\end{equation}
из кинетической нормы устраняет внутреннее противоречие H-M, но не выводит
сам seed. Теперь необходимо независимо построить объёмную вершину и
петлевой проекционный вес
\begin{equation}
 \mathcal J_{\mathrm{req}}
 =\frac{1/3}{|I_\tau|/\pi}
 =5.4027533071\ldots.
\end{equation}
Следовательно, разрыв действительно локализован в вершинной карте, но не
закрыт.

\section{Статус single-vertex утверждения}

Проверенный циклический оператор даёт
\begin{equation}
 \Tr Q_{\mathrm{cycle}}=\pi+\pi^{-1},
\end{equation}
а не \(\pi^2+2\pi\), поэтому именно этот минимальный кандидат проваливается.
Из этого не следует no-go для любой одной градуированной вершины.

Для общего утверждения необходимо заранее определить:
\begin{enumerate}
 \item допустимую алгебру вершин;
 \item разрешённые блоки и градуировки;
 \item единственный ковариантный readout;
 \item запрет на коэффициенты, содержащие целевые числа.
\end{enumerate}
Без этих ограничений один блочно-диагональный оператор может формально
содержать произвольные объёмный и winding-коэффициенты; это target-loaded
запись, но она опровергает возможность универсального no-go без объявления
класса.

\begin{proposition}[Исправленный вердикт H-M]
Гипотеза H-M проходит тест кинематической совместимости:
\begin{equation}
 W=I
 \quad\Longrightarrow\quad
 \left(23+\pi^{-1},\,\frac83\right).
\end{equation}
Она не проходит гейт происхождения меры и не добавляет второго физического
сектора. Перенос сырого seed в вершину является допустимой локализацией
разрыва, но объёмная вершина и \(\mathcal J_{\mathrm{req}}\) остаются
невыведенными. Для \(Q_{\mathrm{cycle}}\) получен частный отрицательный
результат; общий single-vertex no-go не доказан.
\end{proposition}

\begin{protocol}[Следующий вершинный гейт]
Следует заморозить конечномерную алгебру допустимых объёмных и winding-вершин
до сравнения с \(m_\tau\), вывести её из одного действия и проверить, даёт ли
один ковариантный след оба коэффициента без sector-specific проекторов.
Только провал полного объявленного класса позволит сформулировать строгий
single-vertex no-go.
\end{protocol}

Машинный протокол сохранён в
\texttt{s2t\_canonical\_measure\_vertex\_localization\_results.json}.